\documentclass[12pt,a4paper]{article}
\usepackage[utf8]{inputenc}
\usepackage[T2A]{fontenc}
\usepackage[russian]{babel}
\usepackage{amsmath,amssymb,amsthm}
\usepackage{mathtools}
\usepackage{geometry}
\usepackage{hyperref}
\usepackage{cleveref}
\usepackage{enumitem}
\usepackage{tocloft}

\geometry{left=3cm,right=2cm,top=2.5cm,bottom=2.5cm}

\theoremstyle{plain}
\newtheorem{theorem}{Теорема}[section]
\newtheorem{lemma}[theorem]{Лемма}
\newtheorem{proposition}[theorem]{Предложение}
\newtheorem{claim}[theorem]{Утверждение}
\newtheorem{corollary}[theorem]{Следствие}

\theoremstyle{definition}
\newtheorem{definition}[theorem]{Определение}
\newtheorem{remark}[theorem]{Замечание}

\DeclareMathOperator{\Gap}{Gap}
\DeclareMathOperator{\SP}{SP}
\DeclareMathOperator{\Rad}{Rad}
\DeclareMathOperator*{\argmin}{argmin}
\DeclareMathOperator*{\argmax}{argmax}

\newcommand{\E}{\mathbb{E}}
\renewcommand{\Pr}{\mathbb{P}}
\newcommand{\R}{\mathbb{R}}
\newcommand{\cX}{\mathcal{X}}
\newcommand{\cY}{\mathcal{Y}}
\newcommand{\cZ}{\mathcal{Z}}
\newcommand{\cN}{\mathcal{N}}
\newcommand{\cD}{\mathcal{D}}
\newcommand{\hatF}{\hat{F}}
\newcommand{\hatD}{\hat{D}}
\newcommand{\hatG}{\hat{G}}

\begin{document}

% ============================================================
% Титульный лист
% ============================================================
\begin{titlepage}
\centering
\vspace*{2cm}
{\large Санкт-Петербургский государственный университет}\\[0.5cm]
{\large Математико-механический факультет}\\[3cm]
{\LARGE\bfseries Сложность выборки для ERM в стохастических\\[0.3cm]
седловых задачах}\\[1cm]
{\large Обобщение результатов Carmon, Livni, Yehudayoff (2023)\\
на выпукло-вогнутый случай}\\[2cm]
{\large Курсовая работа}\\[3cm]
\begin{flushright}
{\large Выполнил: [Имя студента]}\\[0.3cm]
{\large Научный руководитель: [ФИО]}
\end{flushright}
\vfill
{\large Санкт-Петербург, 2025}
\end{titlepage}

% ============================================================
\tableofcontents
\newpage

% ============================================================
\begin{abstract}
Статья Carmon, Livni, Yehudayoff (2023) решила открытую проблему Фельдмана: для стохастической выпуклой оптимизации (SCO) сложность выборки любого ERM равна $\tilde{O}(d/\varepsilon + 1/\varepsilon^2)$. В настоящей работе мы обобщаем этот результат на стохастические седловые задачи вида $\min_{x\in\cX}\max_{y\in\cY} \E_z[f(x,y,z)]$, где $f$ выпукла по $x$ и вогнута по $y$. Мы показываем, что при $n \ge \tilde{O}((d_x+d_y)/\varepsilon + 1/\varepsilon^2)$ с вероятностью не менее $3/4$ любой Saddle-ERM $(\hat{x},\hat{y})\in\SP(\hatF)$ удовлетворяет $\Gap_F(\hat{x},\hat{y})\le C\varepsilon$ для некоторой абсолютной константы $C>0$. Доказательство следует методологии оригинальной статьи: стохастическое условие первого порядка, пара дивергенций Брегмана, концентрационные неравенства Бернштейна и аргумент $\varepsilon$-сети по совместному пространству $\cX\times\cY$.
\end{abstract}

% ============================================================
\section{Введение}
% ============================================================

\subsection{Контекст}

Стохастическая выпуклая оптимизация (SCO) --- стандартная модель для изучения статистических свойств алгоритмов обучения. В классической постановке задачи минимизации алгоритм ERM (Empirical Risk Minimizer) минимизирует эмпирический риск по выборке, и вопрос состоит в том, как много точек данных необходимо, чтобы любой ERM автоматически хорошо работал на истинном распределении.

Фельдман \cite{Feldman2016} показал, что для ERM в SCO необходимо $\Omega(d/\varepsilon + 1/\varepsilon^2)$ точек, и предложил как открытую проблему вопрос о достаточности этого количества. Carmon, Livni, Yehudayoff \cite{CLY2023} закрыли этот вопрос, доказав, что $\tilde{O}(d/\varepsilon + 1/\varepsilon^2)$ точек \emph{достаточно}. Это разделяет сложность ERM и равномерной сходимости (которая требует $\Theta(d/\varepsilon^2)$) и неожиданно объединяет два режима обучения в одной модели.

\subsection{Наш вклад}

Мы обобщаем результат \cite{CLY2023} на задачу стохастической выпукло-вогнутой оптимизации:
\[
    \min_{x\in\cX}\max_{y\in\cY} F(x,y), \qquad F(x,y) = \E_{z\sim\cD}[f(x,y,z)].
\]
Алгоритм \emph{Saddle-ERM} находит седловую точку эмпирического функционала $\hatF$, а качество решения измеряется через \emph{duality gap}. Основной результат:

\begin{theorem}[Основной результат; точная формулировка --- Теорема~\ref{thm:main}]
\label{thm:intro}
Пусть $\cX\subset\R^{d_x}$, $\cY\subset\R^{d_y}$ --- единичные шары евклидовой нормы, $f_z:\cX\times\cY\to\R$ выпукла по $x$, вогнута по $y$, $1$-липшицева и $|f_z|\le 1$. Для любого $0<\varepsilon<1$, если
\[
    n \ge \tilde{O}\!\left(\frac{d_x+d_y}{\varepsilon}+\frac{1}{\varepsilon^2}\right),
\]
то с вероятностью не менее $3/4$ над $S\sim\cD^n$ любой Saddle-ERM $(\hat{x},\hat{y})\in\SP(\hatF)$ удовлетворяет $\Gap_F(\hat{x},\hat{y})\le C\varepsilon$ для некоторой абсолютной константы $C>0$.
\end{theorem}

Доказательство состоит из семи шагов: (1) постановка задачи и определения, (2) стохастическое условие первого порядка для седловых точек, (3) пара дивергенций Брегмана и их свойства, (4) концентрационная лемма, (5) аналоги ключевых утверждений из \cite{CLY2023}, (6) сборка основной теоремы, (7) обсуждение нижней границы.

% ============================================================
\section{Исходная задача и результат Carmon--Livni--Yehudayoff}
% ============================================================

\subsection{Постановка задачи минимизации}

Пусть $\|\cdot\|$ --- норма на $\R^d$, $K$ --- её единичный шар, $\cZ$ --- конечное пространство данных, $\cD$ --- распределение на $\cZ$. Для каждого $z\in\cZ$ задана функция $f_z:K\to\R$, выпуклая и $L$-липшицева. Популяционный и эмпирический риски:
\[
    F(x) = \E_{z\sim\cD}[f_z(x)], \qquad \hatF(x) = \frac{1}{n}\sum_{i=1}^n f_{z_i}(x).
\]
\emph{ERM} --- любой алгоритм, возвращающий $\hat{x}\in\argmin_{x\in K}\hatF(x)$. Критерий успеха: $F(\hat{x})\le F(x^*)+\varepsilon$.

\subsection{Главный результат}

\begin{theorem}[Carmon, Livni, Yehudayoff \cite{CLY2023}, теорема для $\ell_2$]
\label{thm:CLY}
Пусть $f_z$ выпуклы, $1$-липшицевы по $\|\cdot\|_2$ и $|f_z(x)|\le 1$ для всех $x\in B$. Если
\[
    n \ge n_0 = \frac{3d\ln(40/\varepsilon)}{\varepsilon} + \frac{40}{\varepsilon^2},
\]
то с вероятностью не менее $3/4$ над $S\sim\cD^n$, для всех $x\in B$:
\[
    \hatF(x) \le \hatF(x^*)+\varepsilon \;\Rightarrow\; F(x)\le F(x^*)+40\varepsilon.
\]
\end{theorem}

\subsection{Ключевые идеи доказательства}

Доказательство опирается на три элемента:

\begin{enumerate}
\item \textbf{Стохастическое условие первого порядка.} В минимуме $x^*$ существует субградиент $G\in\partial F(x^*)$ вида $G=\E_z[g_z]$, где $g_z\in\partial f_z(x^*)$, такой что $\langle G,x-x^*\rangle\ge 0$ для всех $x\in K$.

\item \textbf{Дивергенция Брегмана.} $D_F(x,x^*)=F(x)-F(x^*)-\langle G,x-x^*\rangle\ge 0$, ограничена сверху константой. Эмпирическая версия $\hatD(x,x^*)$ является несмещённой оценкой $D_F(x,x^*)$.

\item \textbf{Концентрация.} Если $D_F(x,x^*)\ge 2\varepsilon$, то:
\[
    \Pr\!\left[\hatD(x,x^*)\le\frac{D_F(x,x^*)}{2}\right]\le\exp\!\left(-\frac{D_F(x,x^*)\cdot n}{40}\right).
\]
Ключевой выигрыш: экспонента $\propto D\cdot n$ (а не $D^2\cdot n$), что даёт $d/\varepsilon$ вместо $d/\varepsilon^2$. Затем union bound по $\varepsilon$-сети размера $(1/\varepsilon)^d$ завершает доказательство.
\end{enumerate}

Это разделяет ERM ($\tilde{O}(d/\varepsilon+1/\varepsilon^2)$) и равномерную сходимость ($\Theta(d/\varepsilon^2)$): ERM строго эффективнее.

% ============================================================
\section{Седловая задача: постановка и определения}
% ============================================================

\subsection{Формальная постановка}

Пусть $\|\cdot\|_x$ и $\|\cdot\|_y$ --- нормы на $\R^{d_x}$ и $\R^{d_y}$ соответственно, $\cX$ и $\cY$ --- их единичные шары. Пусть $\cZ$ --- конечное пространство данных, $\cD$ --- распределение на $\cZ$. Для каждого $z\in\cZ$ задана функция $f_z:\cX\times\cY\to\R$, удовлетворяющая:
\begin{itemize}
    \item $f_z(\cdot,y)$ выпукла и $L$-липшицева по $\|\cdot\|_x$ для всех $y$,
    \item $f_z(x,\cdot)$ вогнута и $L$-липшицева по $\|\cdot\|_y$ для всех $x$,
    \item $|f_z(x,y)|\le c$ для всех $(x,y)\in\cX\times\cY$.
\end{itemize}

\begin{definition}[Популяционный функционал]
\[
    F(x,y)=\E_{z\sim\cD}[f_z(x,y)].
\]
Функция $F(\cdot,y)$ выпукла по $x$, $F(x,\cdot)$ вогнута по $y$.
\end{definition}

\begin{definition}[Седловая точка]
Точка $(x^*,y^*)\in\cX\times\cY$ называется \emph{седловой точкой} $F$, если
\[
    F(x^*,y)\le F(x^*,y^*)\le F(x,y^*) \quad\forall x\in\cX,\,y\in\cY.
\]
Множество всех седловых точек обозначим $\SP(F)$.
\end{definition}

\begin{remark}
Седловая точка существует, поскольку $F$ выпукло-вогнута и непрерывна на компактных выпуклых множествах $\cX$, $\cY$ (теорема Сиона--Какутани).
\end{remark}

\begin{definition}[Эмпирический функционал и Saddle-ERM]
По выборке $S=(z_1,\ldots,z_n)\sim\cD^n$ определяется:
\[
    \hatF(x,y)=\frac{1}{n}\sum_{i=1}^n f_{z_i}(x,y).
\]
\emph{Saddle-ERM} --- любой алгоритм, возвращающий $(\hat{x},\hat{y})\in\SP(\hatF)$, то есть седловую точку эмпирического функционала. Такая точка существует по тем же причинам, что и для $F$.
\end{definition}

\begin{definition}[Duality gap]
\label{def:gap}
Для точки $(x,y)\in\cX\times\cY$ определим:
\[
    \Gap_F(x,y)=\max_{y'\in\cY}F(x,y')-\min_{x'\in\cX}F(x',y)\ge 0.
\]
$\Gap_F(x,y)=0$ тогда и только тогда, когда $(x,y)\in\SP(F)$.
\emph{Saddle-ERM успешен} с параметром $\varepsilon>0$, если $\Gap_F(\hat{x},\hat{y})\le\varepsilon$.
\end{definition}

\begin{remark}[Связь с задачей минимизации]
Используя $(x^*,y^*)\in\SP(F)$:
\begin{align}
    \Gap_F(x,y)
    &\le \bigl[F(x,y^*)-F(x^*,y^*)\bigr]+\bigl[F(x^*,y^*)-F(x^*,y)\bigr].
    \label{eq:gap_decomp}
\end{align}
Каждое слагаемое неотрицательно и характеризует отклонение от $x^*$ и от $y^*$ соответственно.
\end{remark}

% ============================================================
\section{Вспомогательные результаты}
% ============================================================

\subsection{Субградиентное исчисление}

Субградиент выпуклой функции $h:\cX\to\R$ в точке $x_0$:
\[
    \partial_x h(x_0)=\{g\in\R^{d_x}:h(x)\ge h(x_0)+\langle g,x-x_0\rangle\;\forall x\in\cX\}.
\]
Для вогнутой функции $h:\cY\to\R$ субградиент определяем через $-h$:
\[
    \partial_y^- h(y_0)=\{g\in\R^{d_y}:h(y)\le h(y_0)+\langle g,y-y_0\rangle\;\forall y\in\cY\}.
\]
Если $h$ является $L$-липшицевой по $\|\cdot\|_x$, то $\|g\|_{x,*}\le L$ для всех $g\in\partial_x h(x_0)$, где $\|\cdot\|_{x,*}$ --- двойственная норма.

\subsection{Стохастическое условие первого порядка}

\begin{proposition}[Стохастическое условие первого порядка для седловой точки]
\label{prop:firstorder}
Пусть $(x^*,y^*)\in\SP(F)$. Тогда для каждого $z\in\cZ$ существуют
\[
    g_z^x\in\partial_x f_z(x^*,y^*),\qquad g_z^y\in\partial_y^- f_z(x^*,y^*)
\]
с $\|g_z^x\|_{x,*}\le L$, $\|g_z^y\|_{y,*}\le L$, такие что, обозначив
$G_x=\E_z[g_z^x]$, $G_y=\E_z[g_z^y]$, выполняется:
\begin{align}
    \langle G_x,x-x^*\rangle&\ge 0\quad\forall x\in\cX, \label{eq:foc_x}\\
    \langle G_y,y-y^*\rangle&\ge 0\quad\forall y\in\cY. \label{eq:foc_y}
\end{align}
\end{proposition}

\begin{proof}
\textbf{Условие \eqref{eq:foc_x}.} Поскольку $(x^*,y^*)$ --- седловая точка, $x^*$ минимизирует $F(\cdot,y^*)$ на $\cX$. По условию первого порядка для выпуклой задачи (см. \cite{Rockafellar1997}) существует $G_x\in\partial_x F(x^*,y^*)$ такой, что $\langle G_x,x-x^*\rangle\ge 0$ для всех $x\in\cX$. По теореме о сумме субградиентов \cite{Rockafellar1997}, применённой к $F(\cdot,y^*)=\E_z[f_z(\cdot,y^*)]$, для каждого $z$ существует $g_z^x\in\partial_x f_z(x^*,y^*)$ такой, что $G_x=\E_z[g_z^x]$. Липшицевость даёт $\|g_z^x\|_{x,*}\le L$.

\textbf{Условие \eqref{eq:foc_y}.} Аналогично: $y^*$ максимизирует $F(x^*,\cdot)$ на $\cY$, то есть является минимизатором $-F(x^*,\cdot)$. Применяем условие первого порядка к выпуклой функции $-F(x^*,\cdot)$: существует $G_y\in\partial_y(-F)(x^*,y^*)=\partial_y^- F(x^*,y^*)$ такой, что $\langle G_y,y-y^*\rangle\ge 0$. По теореме о сумме субградиентов, $G_y=\E_z[g_z^y]$, $g_z^y\in\partial_y(-f_z)(x^*,y^*)$.
\end{proof}

\begin{remark}
Условия \eqref{eq:foc_x}--\eqref{eq:foc_y} образуют \emph{вариационное неравенство}, характеризующее седловую точку. Это двумерный аналог условия $\langle G,x-x^*\rangle\ge 0$ из \cite{CLY2023}.
\end{remark}

\subsection{Пара дивергенций Брегмана}

Зафиксируем $(x^*,y^*)\in\SP(F)$ и субградиенты $G_x,G_y$ из Предложения~\ref{prop:firstorder}. Зафиксируем также выборочные версии $\hatG_x=\frac{1}{n}\sum_i g_{z_i}^x$, $\hatG_y=\frac{1}{n}\sum_i g_{z_i}^y$.

\begin{definition}[Пара дивергенций Брегмана]
\label{def:bregman}
\begin{align*}
    D^x(x,x^*)&=F(x,y^*)-F(x^*,y^*)-\langle G_x,x-x^*\rangle,\\
    D^y(y,y^*)&=F(x^*,y^*)-F(x^*,y)-\langle G_y,y-y^*\rangle.
\end{align*}
\emph{Поточечные} (стохастические) версии:
\begin{align*}
    D_z^x(x,x^*)&=f_z(x,y^*)-f_z(x^*,y^*)-\langle g_z^x,x-x^*\rangle,\\
    D_z^y(y,y^*)&=f_z(x^*,y^*)-f_z(x^*,y)-\langle g_z^y,y-y^*\rangle.
\end{align*}
\emph{Эмпирические} версии:
\begin{align*}
    \hatD^x(x,x^*)&=\hatF(x,y^*)-\hatF(x^*,y^*)-\langle\hatG_x,x-x^*\rangle,\\
    \hatD^y(y,y^*)&=\hatF(x^*,y^*)-\hatF(x^*,y)-\langle\hatG_y,y-y^*\rangle.
\end{align*}
\end{definition}

\begin{lemma}[Неотрицательность]
\label{lem:nonneg}
Для всех $(x,y)\in\cX\times\cY$:
\[
    D^x(x,x^*)\ge 0,\quad D^y(y,y^*)\ge 0,\quad D_z^x\ge 0,\quad D_z^y\ge 0,\quad \hatD^x\ge 0,\quad \hatD^y\ge 0.
\]
\end{lemma}

\begin{proof}
$D^x\ge 0$: функция $F(\cdot,y^*)$ выпукла, $G_x\in\partial_x F(x^*,y^*)$, поэтому по определению субградиента $F(x,y^*)\ge F(x^*,y^*)+\langle G_x,x-x^*\rangle$.

$D^y\ge 0$: функция $-F(x^*,\cdot)$ выпукла, $G_y\in\partial_y(-F)(x^*,y^*)$, поэтому $-F(x^*,y)\ge-F(x^*,y^*)+\langle G_y,y-y^*\rangle$, откуда $D^y\ge 0$.

Для $D_z^x\ge 0$ и $D_z^y\ge 0$ рассуждение аналогично для $f_z$ вместо $F$. Неотрицательность $\hatD^x,\hatD^y$ следует как среднее неотрицательных слагаемых.
\end{proof}

\begin{lemma}[Ограниченность]
\label{lem:bounded}
Для всех $z\in\cZ$ и $(x,y)\in\cX\times\cY$:
\[
    0\le D_z^x(x,x^*)\le 4c,\qquad 0\le D_z^y(y,y^*)\le 4c.
\]
Следовательно, $0\le D^x(x,x^*)\le 4c$ и $0\le D^y(y,y^*)\le 4c$.
\end{lemma}

\begin{proof}
Оценим $D_z^x$ сверху. Так как $|f_z|\le c$:
\[
    f_z(x,y^*)-f_z(x^*,y^*)\le 2c.
\]
Из липшицевости $f_z(\cdot,y^*)$ при $x,x^*\in\cX$:
\[
    |\langle g_z^x,x-x^*\rangle|\le\|g_z^x\|_{x,*}\cdot\|x-x^*\|_x\le L\cdot 2=2L.
\]
В случае нормированных шаров ($L=c=1$) получаем $D_z^x\le 2+2=4$, откуда в общем случае $D_z^x\le 4c$. Аналогично для $D_z^y$.
\end{proof}

\begin{corollary}[Несмещённость]
\label{cor:unbiased}
$\E_{S\sim\cD^n}[\hatD^x(x,x^*)]=D^x(x,x^*)$ и $\E_{S\sim\cD^n}[\hatD^y(y,y^*)]=D^y(y,y^*)$.
\end{corollary}

\begin{proof}
$\hatD^x(x,x^*)=\frac{1}{n}\sum_i D_{z_i}^x(x,x^*)$, поскольку $\hatF(x,y^*)-\hatF(x^*,y^*)-\langle\hatG_x,x-x^*\rangle=\frac{1}{n}\sum_i[f_{z_i}(x,y^*)-f_{z_i}(x^*,y^*)-\langle g_{z_i}^x,x-x^*\rangle]$. Так как $z_i$ i.i.d. из $\cD$, то $\E[\hatD^x]=\E_z[D_z^x]=D^x$.
\end{proof}

\subsection{Связь дивергенций с duality gap}

\begin{lemma}[Разложение gap через дивергенции]
\label{lem:gap_bregman}
Для любой точки $(x,y)\in\cX\times\cY$:
\[
    \Gap_F(x,y)\le D^x(x,x^*)+D^y(y,y^*)+\langle G_x,x-x^*\rangle+\langle G_y,y-y^*\rangle.
\]
\end{lemma}

\begin{proof}
Из \eqref{eq:gap_decomp}:
\begin{align*}
    \Gap_F(x,y)
    &\le [F(x,y^*)-F(x^*,y^*)]+[F(x^*,y^*)-F(x^*,y)]\\
    &=[D^x(x,x^*)+\langle G_x,x-x^*\rangle]+[D^y(y,y^*)+\langle G_y,y-y^*\rangle].\qedhere
\end{align*}
\end{proof}

\subsection{Концентрационная лемма}

\begin{claim}[Концентрация дивергенций Брегмана]
\label{claim:concentration}
При условиях раздела~3 для любых $(x,y)\in\cX\times\cY$:
\begin{align}
    \Pr\!\left[\hatD^x(x,x^*)\le\frac{D^x(x,x^*)}{2}\right]&\le\exp\!\left(-\frac{D^x(x,x^*)\cdot n}{40c}\right),\label{eq:conc_x}\\
    \Pr\!\left[\hatD^y(y,y^*)\le\frac{D^y(y,y^*)}{2}\right]&\le\exp\!\left(-\frac{D^y(y,y^*)\cdot n}{40c}\right).\label{eq:conc_y}
\end{align}
\end{claim}

\begin{proof}
Докажем \eqref{eq:conc_x}; доказательство \eqref{eq:conc_y} симметрично.

Запишем $\hatD^x(x,x^*)=\frac{1}{n}\sum_{j=1}^n R_j$, где $R_j=D_{z_j}^x(x,x^*)$ --- i.i.d. случайные величины. По лемме~\ref{lem:nonneg} и~\ref{lem:bounded}: $R_j\in[0,4c]$ и $\E[R_j]=\lambda:=D^x(x,x^*)$.

По неравенству Бхатиа--Дэвиса \cite{BhatiaDavis2000}: $\mathrm{Var}(R_j)\le(4c-\lambda)\lambda\le 4c\lambda$.

По неравенству Бернштейна:
\begin{align*}
    \Pr\!\left[\hatD^x(x,x^*)\le\frac{\lambda}{2}\right]
    &=\Pr\!\left[\sum_j(R_j-\lambda)\le-\frac{\lambda n}{2}\right]\\
    &\le\exp\!\left(-\frac{(\lambda n/2)^2}{2\cdot 4c\lambda n+2\cdot 4c\cdot(\lambda n/2)/3}\right)\\
    &=\exp\!\left(-\frac{\lambda^2 n^2/4}{8c\lambda n+4c\lambda n/3}\right)
    \le\exp\!\left(-\frac{\lambda n}{40c}\right).\qedhere
\end{align*}
\end{proof}

% ============================================================
\section{Аналоги ключевых утверждений для седловой задачи}
% ============================================================

В этом разделе мы доказываем аналоги Claims~5,\,6 из \cite{CLY2023} для седлового случая.

\subsection{Контроль линейного члена}

Обозначим $\delta_x=\|\hatG_x-G_x\|_{x,*}$ и $\delta_y=\|\hatG_y-G_y\|_{y,*}$. Определим \emph{хорошее событие}:
\[
    H = \{\delta_x\le\varepsilon\}\cap\{\delta_y\le\varepsilon\}.
\]

\begin{claim}[Концентрация эмпирических субградиентов]
\label{claim:grad_conc}
\[
    \Pr[\delta_x^2>\varepsilon^2]\le\frac{4L^2}{\varepsilon^2 n},\qquad
    \Pr[\delta_y^2>\varepsilon^2]\le\frac{4L^2}{\varepsilon^2 n}.
\]
В частности, $\Pr[\neg H]\le 8L^2/(\varepsilon^2 n)$.
\end{claim}

\begin{proof}
$\E\|\hatG_x-G_x\|_{x,*}^2=\frac{1}{n^2}\sum_j\E\|g_{z_j}^x-G_x\|_{x,*}^2\le\frac{4L^2}{n}$ (так как $\|g_z^x\|_{x,*}\le L$). По неравенству Маркова: $\Pr[\delta_x^2>\varepsilon^2]\le\frac{4L^2}{\varepsilon^2 n}$. Аналогично для $\delta_y$.
\end{proof}

\subsection{Утверждение о малости эмпирических дивергенций (аналог Claim~6)}

\begin{claim}[Saddle-ERM имеет малые эмпирические дивергенции]
\label{claim:saddle_claim6}
Пусть $(\hat{x},\hat{y})\in\SP(\hatF)$. Тогда на событии $H$:
\[
    \hatD^x(\hat{x},x^*)+\hatD^y(\hat{y},y^*)\le 4\varepsilon.
\]
\end{claim}

\begin{proof}
Заметим:
\begin{align*}
    \hatD^x(\hat{x},x^*)+\hatD^y(\hat{y},y^*)
    &=\hatF(\hat{x},y^*)-\hatF(x^*,y^*)-\langle\hatG_x,\hat{x}-x^*\rangle\\
    &\quad+\hatF(x^*,y^*)-\hatF(x^*,\hat{y})-\langle\hatG_y,\hat{y}-y^*\rangle\\
    &=\hatF(\hat{x},y^*)-\hatF(x^*,\hat{y})-\langle\hatG_x,\hat{x}-x^*\rangle-\langle\hatG_y,\hat{y}-y^*\rangle.
\end{align*}

Из седловой точки $(\hat{x},\hat{y})\in\SP(\hatF)$: $\hatF(\hat{x},y^*)\le\hatF(\hat{x},\hat{y})\le\hatF(x^*,\hat{y})$, значит $\hatF(\hat{x},y^*)-\hatF(x^*,\hat{y})\le 0$.

Поэтому:
\begin{align*}
    \hatD^x+\hatD^y
    &\le-\langle\hatG_x,\hat{x}-x^*\rangle-\langle\hatG_y,\hat{y}-y^*\rangle\\
    &=-\langle\hatG_x-G_x,\hat{x}-x^*\rangle-\langle G_x,\hat{x}-x^*\rangle-\langle\hatG_y-G_y,\hat{y}-y^*\rangle-\langle G_y,\hat{y}-y^*\rangle.
\end{align*}

Из условий \eqref{eq:foc_x}--\eqref{eq:foc_y}: $-\langle G_x,\hat{x}-x^*\rangle\le 0$ и $-\langle G_y,\hat{y}-y^*\rangle\le 0$. На событии $H$:
\[
    |\langle\hatG_x-G_x,\hat{x}-x^*\rangle|\le\|\hatG_x-G_x\|_{x,*}\cdot\|\hat{x}-x^*\|_x\le\varepsilon\cdot 2=2\varepsilon,
\]
аналогично $|\langle\hatG_y-G_y,\hat{y}-y^*\rangle|\le 2\varepsilon$.

Итого: $\hatD^x(\hat{x},x^*)+\hatD^y(\hat{y},y^*)\le 2\varepsilon+0+2\varepsilon+0=4\varepsilon$.
\end{proof}

\subsection{Утверждение о связи gap с дивергенциями (аналог Claim~5)}

\begin{claim}[Большой gap влечёт большие дивергенции]
\label{claim:saddle_claim5}
Пусть $(\hat{x},\hat{y})\in\SP(\hatF)$. На событии $H$:
\[
    \Gap_F(\hat{x},\hat{y})\le D^x(\hat{x},x^*)+D^y(\hat{y},y^*)+4\varepsilon.
\]
В частности, если $D^x(\hat{x},x^*)+D^y(\hat{y},y^*)\le C\varepsilon$, то $\Gap_F(\hat{x},\hat{y})\le(C+4)\varepsilon$.
\end{claim}

\begin{proof}
Из леммы~\ref{lem:gap_bregman}:
\[
    \Gap_F(\hat{x},\hat{y})\le D^x+D^y+\langle G_x,\hat{x}-x^*\rangle+\langle G_y,\hat{y}-y^*\rangle.
\]

Оценим линейные члены. Из определения $\hatD^x$:
\[
    \langle\hatG_x,\hat{x}-x^*\rangle=\hatF(\hat{x},y^*)-\hatF(x^*,y^*)-\hatD^x(\hat{x},x^*).
\]
Аналогично:
\[
    \langle\hatG_y,\hat{y}-y^*\rangle=\hatF(x^*,y^*)-\hatF(x^*,\hat{y})-\hatD^y(\hat{y},y^*).
\]
Суммируя:
\[
    \langle\hatG_x,\hat{x}-x^*\rangle+\langle\hatG_y,\hat{y}-y^*\rangle=\hatF(\hat{x},y^*)-\hatF(x^*,\hat{y})-\hatD^x-\hatD^y\le 0
\]
(используем $\hatF(\hat{x},y^*)\le\hatF(x^*,\hat{y})$ из седловой точки и $\hatD^x,\hatD^y\ge 0$).

Теперь:
\begin{align*}
    \langle G_x,\hat{x}-x^*\rangle+\langle G_y,\hat{y}-y^*\rangle
    &=\langle G_x-\hatG_x,\hat{x}-x^*\rangle+\langle\hatG_x,\hat{x}-x^*\rangle\\
    &\quad+\langle G_y-\hatG_y,\hat{y}-y^*\rangle+\langle\hatG_y,\hat{y}-y^*\rangle\\
    &\le 2\delta_x+0+2\delta_y+0\le 4\varepsilon
\end{align*}
на событии $H$. Подставляя, получаем утверждение.
\end{proof}

% ============================================================
\section{Основная теорема}
% ============================================================

\subsection{Формулировка}

\begin{theorem}[Сложность выборки Saddle-ERM, случай $\ell_2$]
\label{thm:main}
Пусть $\|\cdot\|_x=\|\cdot\|_y=\|\cdot\|_2$, $L=1$, $c=1$. Для любого $0<\varepsilon<1$ если
\[
    n\ge n_0=\frac{3(d_x+d_y)\ln(40/\varepsilon)}{\varepsilon}+\frac{80}{\varepsilon^2},
\]
то с вероятностью не менее $3/4$ над $S\sim\cD^n$ любой Saddle-ERM $(\hat{x},\hat{y})\in\SP(\hatF)$ удовлетворяет:
\[
    \Gap_F(\hat{x},\hat{y})\le 50\varepsilon.
\]
\end{theorem}

\subsection{Доказательство}

\begin{proof}
Определим \emph{плохое событие}:
\[
    E=\bigl\{(\hat{x},\hat{y})\in\SP(\hatF):\Gap_F(\hat{x},\hat{y})>50\varepsilon\bigr\}.
\]
Покажем, что $\Pr[E]\le 1/4$.

\textbf{Шаг 1: Хорошее событие $H$.}
По утверждению~\ref{claim:grad_conc}:
\[
    \Pr[\neg H]\le\frac{8}{\varepsilon^2 n}\le\frac{1}{8}
\]
при $n\ge 64/\varepsilon^2$.

\textbf{Шаг 2: Работаем на событии $H$.}
Предположим, что $E\cap H$ произошло. По утверждению~\ref{claim:saddle_claim5}: если $\Gap_F(\hat{x},\hat{y})>50\varepsilon$, то
\[
    D^x(\hat{x},x^*)+D^y(\hat{y},y^*)>46\varepsilon.
\]

\textbf{Шаг 3: $\varepsilon$-сеть по $\cX\times\cY$.}
Пусть $\cN_x$ --- $(\varepsilon/3)$-сеть по $\cX$ и $\cN_y$ --- $(\varepsilon/3)$-сеть по $\cY$ (обе по $\ell_2$). По лемме о покрывающих числах \cite{CLY2023}:
\[
    |\cN_x|\le\left(\frac{12}{\varepsilon}\right)^{d_x},\qquad
    |\cN_y|\le\left(\frac{12}{\varepsilon}\right)^{d_y}.
\]
Совместная сеть $\cN=\cN_x\times\cN_y$ имеет размер $|\cN|\le(12/\varepsilon)^{d_x+d_y}$.

Поскольку $F(\cdot,y^*)$ и $F(x^*,\cdot)$ являются $1$-липшицевыми (по $\ell_2$), а $\hatF(\cdot,y^*)$ и $\hatF(x^*,\cdot)$ --- тоже, дивергенции $D^x(\cdot,x^*)$, $D^y(\cdot,y^*)$, $\hatD^x(\cdot,x^*)$, $\hatD^y(\cdot,y^*)$ являются $2$-липшицевыми по $\ell_2$ (оценка из ограниченности субградиентов). Значит, для $\|\hat{x}-\tilde{x}\|_2\le\varepsilon/3$:
\[
    |D^x(\hat{x},x^*)-D^x(\tilde{x},x^*)|\le\frac{2\varepsilon}{3},
\]
и аналогично для остальных дивергенций.

Для сетевой точки $(\tilde{x},\tilde{y})\in\cN$ с $\|\hat{x}-\tilde{x}\|_2\le\varepsilon/3$, $\|\hat{y}-\tilde{y}\|_2\le\varepsilon/3$:
\[
    D^x(\tilde{x},x^*)+D^y(\tilde{y},y^*)\ge 46\varepsilon-\frac{4\varepsilon}{3}>44\varepsilon,
\]
\[
    \hatD^x(\tilde{x},x^*)+\hatD^y(\tilde{y},y^*)\le 4\varepsilon+\frac{4\varepsilon}{3}<6\varepsilon.
\]
(Последнее --- из утверждения~\ref{claim:saddle_claim6} и липшицевости $\hatD^x,\hatD^y$.)

\textbf{Шаг 4: Концентрационный аргумент.}
Из $D^x(\tilde{x},x^*)+D^y(\tilde{y},y^*)>44\varepsilon$ следует: хотя бы одно из $D^x(\tilde{x},x^*)>22\varepsilon$ или $D^y(\tilde{y},y^*)>22\varepsilon$. Предположим $D^x(\tilde{x},x^*)>22\varepsilon$ (случай $D^y>22\varepsilon$ симметричен).

По утверждению~\ref{claim:concentration}:
\[
    \Pr\!\left[\hatD^x(\tilde{x},x^*)\le\frac{D^x(\tilde{x},x^*)}{2}\right]\le\exp\!\left(-\frac{22\varepsilon\cdot n}{40}\right)=\exp\!\left(-\frac{11\varepsilon n}{20}\right).
\]
Но мы показали, что на $E\cap H$ выполнено $\hatD^x(\tilde{x},x^*)\le 6\varepsilon<11\varepsilon=D^x/2$. Значит, событие $E\cap H$ влечёт для соответствующего $(\tilde{x},\tilde{y})\in\cN$: $\hatD^x(\tilde{x},x^*)\le D^x(\tilde{x},x^*)/2$.

Обозначим: $E_{(\tilde{x},\tilde{y})}$ --- событие $\{\hatD^x(\tilde{x},x^*)\le D^x(\tilde{x},x^*)/2\}$ или $\{\hatD^y(\tilde{y},y^*)\le D^y(\tilde{y},y^*)/2\}$ (в зависимости от того, какая дивергенция велика). Тогда:
\[
    E\cap H\subseteq\bigcup_{(\tilde{x},\tilde{y})\in\cN}E_{(\tilde{x},\tilde{y})}.
\]

\textbf{Шаг 5: Union bound.}
По union bound:
\[
    \Pr[E\cap H]\le|\cN|\cdot\exp\!\left(-\frac{11\varepsilon n}{20}\right)\le\left(\frac{12}{\varepsilon}\right)^{d_x+d_y}\exp\!\left(-\frac{11\varepsilon n}{20}\right).
\]
Для $n\ge\frac{3(d_x+d_y)\ln(40/\varepsilon)}{\varepsilon}$ правая часть $\le 1/8$.

\textbf{Итог:}
\[
    \Pr[E]\le\Pr[\neg H]+\Pr[E\cap H]\le\frac{1}{8}+\frac{1}{8}=\frac{1}{4}.\qedhere
\]
\end{proof}

\subsection{Обобщение на произвольные нормы}

\begin{theorem}[Общие нормы]
\label{thm:general_norm}
Пусть $\|\cdot\|_x$, $\|\cdot\|_y$ --- произвольные нормы. Для любых $0<\varepsilon,\delta<1$, если
\[
    n\ge n_0=\frac{12cd_x}{\varepsilon}\ln\frac{3L}{\varepsilon}+\frac{12cd_y}{\varepsilon}\ln\frac{3L}{\varepsilon}+\Rad_{\cX}^{-1}\!\left(\frac{\varepsilon}{2L}\right)+\Rad_{\cY}^{-1}\!\left(\frac{\varepsilon}{2L}\right)+\frac{16c^2}{\varepsilon^2}\ln\frac{8}{\delta},
\]
то с вероятностью не менее $1-\delta$ любой Saddle-ERM $(\hat{x},\hat{y})$ удовлетворяет $\Gap_F(\hat{x},\hat{y})\le C\varepsilon$ для абсолютной константы $C$.
\end{theorem}

\begin{proof}[Набросок доказательства]
Вместо $\|\hatG_x-G_x\|_{x,*}\le\varepsilon$ и $\|\hatG_y-G_y\|_{y,*}\le\varepsilon$ (которые контролировались через дисперсию в случае $\ell_2$) используется понятие \emph{representativeness} из \cite{CLY2023}:
\[
    \mathrm{Rep}_x(S)=\sup_{x\in\cX}(\mathcal{L}_x(x)-\hat{\mathcal{L}}_x(x)),\quad
    \mathrm{Rep}_y(S)=\sup_{y\in\cY}(\mathcal{L}_y(y)-\hat{\mathcal{L}}_y(y)),
\]
где $\mathcal{L}_x$, $\mathcal{L}_y$ --- линейные аппроксимации функционалов, связанные с субградиентами $g_z^x$, $g_z^y$ соответственно. По следствию~1 из \cite{CLY2023} (применённому дважды, к каждой компоненте):
\[
    \mathrm{Rep}_x(S)\le 2L\cdot\Rad(\cX,n)+c\sqrt{\frac{2\ln(4/\delta)}{n}}
\]
и аналогично для $y$-компоненты. Это заменяет контроль $\hatG_x,\hatG_y$ в доказательстве теоремы~\ref{thm:main}. Вместо усечённых дивергенций Брегмана используются усечённые версии $T^x$, $T^y$ (аналог раздела~6 из \cite{CLY2023}). Все остальные шаги переносятся дословно.
\end{proof}

% ============================================================
\section{Нижняя граница и разделение с равномерной сходимостью}
% ============================================================

\subsection{Нижняя граница}

\begin{theorem}[Нижняя граница для Saddle-ERM, неформально]
\label{thm:lower}
Для любого алгоритма Saddle-ERM существует распределение $\cD$ на классе выпукло-вогнутых функций на $\R^{d_x}\times\R^{d_y}$ такой, что при
\[
    n\le c_0\cdot\frac{d_x+d_y}{\varepsilon}
\]
с вероятностью не менее $1/2$: $\Gap_F(\hat{x},\hat{y})>\varepsilon$.
\end{theorem}

\begin{proof}[Конструкция]
Адаптируем конструкцию Фельдмана \cite{Feldman2016} из задачи минимизации. Там задача $\min_{x\in B^d} F(x)$ сводилась к различению $d$ задач на координатных гиперплоскостях. Для седловой задачи берём функцию $f_z(x,y)=f_z^{(x)}(x)-f_z^{(y)}(y)$, где $f_z^{(x)}$ --- конструкция Фельдмана по $x$, $f_z^{(y)}$ --- по $y$, независимые. Тогда:
\begin{itemize}
    \item $f_z(x,y)$ выпукла по $x$ и вогнута по $y$,
    \item любой Saddle-ERM должен одновременно хорошо приближать $\argmin_x F(x,y^*)$ и $\argmax_y F(x^*,y)$,
    \item поскольку задачи по $x$ и по $y$ независимы, нижняя граница на число точек для каждой из них прибавляется, откуда $n=\Omega((d_x+d_y)/\varepsilon)$.
\end{itemize}
Член $1/\varepsilon^2$ наследуется из нижней границы для агностического обучения, применённой к любой из компонент.
\end{proof}

\begin{remark}
Строгое доказательство нижней границы требует точного анализа того, что информация о $x$-компоненте не помогает оценить $y$-компоненту (и наоборот) при независимой конструкции. Это следует из независимости выборок, участвующих в двух подзадачах, что выполняется при нашей конструкции.
\end{remark}

\subsection{Разделение с равномерной сходимостью}

\begin{theorem}[Разделение Saddle-ERM и равномерной сходимости]
\label{thm:separation}
Существуют классы выпукло-вогнутых функций такие, что:
\begin{enumerate}
    \item Любой Saddle-ERM успешен при $n=\tilde{O}((d_x+d_y)/\varepsilon+1/\varepsilon^2)$.
    \item Равномерная сходимость требует $n=\Omega((d_x+d_y)/\varepsilon^2)$.
\end{enumerate}
В частности, Saddle-ERM строго эффективнее равномерной сходимости.
\end{theorem}

\begin{proof}[Набросок]
Верхняя граница для Saddle-ERM --- теорема~\ref{thm:main}. Нижняя граница для равномерной сходимости следует из того, что равномерная сходимость $\sup_{(x,y)\in\cX\times\cY}|F(x,y)-\hatF(x,y)|\le\varepsilon$ требует $n=\Omega((d_x+d_y)/\varepsilon^2)$ по размерностным аргументам (аналог теоремы Фельдмана \cite{Feldman2016}). Действительно, для разделения равномерной сходимости и Saddle-ERM мы можем использовать ту же конструкцию с аддитивной функцией: равномерная сходимость должна контролировать все $(x,y)$-пары одновременно, что требует в $1/\varepsilon$ раз больше точек по $d$-размерному члену.
\end{proof}

% ============================================================
\section{Заключение}
% ============================================================

В работе мы обобщили результат Carmon, Livni, Yehudayoff \cite{CLY2023} на стохастические седловые задачи.

\textbf{Основные результаты:}
\begin{enumerate}
    \item Определены понятия Saddle-ERM и duality gap как критерия успеха в стохастической выпукло-вогнутой оптимизации.
    \item Доказано стохастическое условие первого порядка для седловых точек (Предложение~\ref{prop:firstorder}), обобщающее Proposition~3 из \cite{CLY2023}.
    \item Введена пара дивергенций Брегмана $(D^x,D^y)$; доказаны их неотрицательность (лемма~\ref{lem:nonneg}), ограниченность (лемма~\ref{lem:bounded}) и несмещённость (следствие~\ref{cor:unbiased}).
    \item Доказана концентрационная лемма (утверждение~\ref{claim:concentration}) --- аналог Claim~7 из \cite{CLY2023}.
    \item Доказаны аналоги Claims~5,\,6: утверждения~\ref{claim:saddle_claim6}--\ref{claim:saddle_claim5}, устанавливающие связь между Saddle-ERM и дивергенциями.
    \item Доказана основная теорема~\ref{thm:main}: сложность выборки Saddle-ERM равна $\tilde{O}((d_x+d_y)/\varepsilon+1/\varepsilon^2)$.
    \item Обсуждены нижняя граница и разделение с равномерной сходимостью.
\end{enumerate}

\textbf{Главный структурный вывод.} Ключевой шаг оригинальной статьи --- доказательство того, что $\hat{D}$ мало для ERM, --- переносится на седловой случай через тождество:
\[
    \hatD^x(\hat{x},x^*)+\hatD^y(\hat{y},y^*)=\hatF(\hat{x},y^*)-\hatF(x^*,\hat{y})-\langle\hatG_x,\hat{x}-x^*\rangle-\langle\hatG_y,\hat{y}-y^*\rangle,
\]
где первое слагаемое $\le 0$ по седловому условию, а линейные члены малы по концентрации. Это позволяет применить стандартный аргумент $\varepsilon$-сети по совместному пространству $\cX\times\cY$.

\textbf{Открытые вопросы.} Строгое доказательство нижней границы $\Omega((d_x+d_y)/\varepsilon)$, точный аналог теоремы~\ref{thm:general_norm} с проверкой совместной Радемахеровской сложности, а также вопрос о разделении с равномерной сходимостью в содержательном классе функций остаются направлениями дальнейших исследований.

% ============================================================
\begin{thebibliography}{99}

\bibitem{CLY2023}
D.~Carmon, R.~Livni, A.~Yehudayoff.
\emph{The Sample Complexity of ERMs in Stochastic Convex Optimization}.
arXiv:2311.05398, 2023.

\bibitem{Feldman2016}
V.~Feldman.
\emph{Generalization of ERM in Stochastic Convex Optimization: The Dimension Strikes Back}.
Advances in Neural Information Processing Systems 29, 2016.

\bibitem{Rockafellar1997}
R.~T.~Rockafellar.
\emph{Convex Analysis}.
Princeton University Press, 1997.

\bibitem{BhatiaDavis2000}
R.~Bhatia, C.~Davis.
\emph{A better bound on the variance}.
American Mathematical Monthly, 107(4):353--357, 2000.

\bibitem{ShalevBenDavid2014}
S.~Shalev-Shwartz, S.~Ben-David.
\emph{Understanding Machine Learning: From Theory to Algorithms}.
Cambridge University Press, 2014.

\bibitem{Bubeck2015}
S.~Bubeck.
\emph{Convex Optimization: Algorithms and Complexity}.
Foundations and Trends in Machine Learning, 8(3--4):231--357, 2015.

\bibitem{Nemirovski2009}
A.~Nemirovski, A.~Juditsky, G.~Lan, A.~Shapiro.
\emph{Robust Stochastic Approximation Approach to Stochastic Programming}.
SIAM Journal on Optimization, 19(4):1574--1609, 2009.

\end{thebibliography}

\end{document}
